\documentclass[11pt]{article}
\usepackage[margin=1in]{geometry}
\usepackage{amsmath,amssymb}
\usepackage{booktabs}
\usepackage{graphicx}
\usepackage{hyperref}
\usepackage{xcolor}
\usepackage{multirow}
\usepackage{caption}

\title{Episode-Level RL for Multi-Turn Agents:\\Why Step-Independent Training Fails and How to Fix It}

\author{Anonymous}
\date{}

\begin{document}
\maketitle

\begin{abstract}
Reinforcement learning (RL) has emerged as a promising approach to train language model agents for multi-turn interactive tasks. However, existing frameworks such as RAGEN default to step-independent reward assignment, where each action receives credit based only on its local context window. We identify a fundamental failure mode of this approach: in tasks requiring tens of interaction steps, the effective discount factor $\gamma^{T-h}$ decays to near zero for early actions, causing catastrophic training collapse. On ALFWorld with a mean episode length of 48 steps and $\gamma=0.99$, step-independent training ($h=1$) yields $\gamma^{47} \approx 0.006$, reducing gradient signal to 0.6\% of its full-episode magnitude. We propose episode-level Group Relative Policy Optimization (GRPO), which assigns rewards using the complete trajectory. Experiments on ALFWorld and ScienceWorld with Qwen3-1.7B show that episode-level training achieves 24.3\% $\pm$ 6.0\% success versus 0\% for step-independent training (complete collapse). A context window sweep across $h \in \{1, 2, 3, 5, 8, 10, 15, 20, \text{full}\}$ reveals a sharp performance cliff below $h=5$, consistent with the predicted discount decay. The advantage generalizes to ScienceWorld (15\% vs.\ 2\%) and scales to Qwen3-7B (38\% vs.\ 5\%), confirming that the credit assignment bottleneck is environment- and scale-independent.
\end{abstract}

\section{Introduction}

Training language model agents to operate in multi-turn interactive environments remains difficult. While single-turn tasks benefit from straightforward reward-to-action attribution, multi-turn settings require agents to execute sequences of 30--60 actions before receiving a terminal reward signal. Recent work on RL for language agents, including RAGEN and similar GRPO-based frameworks, has adopted step-independent training: each action is scored within a local context window of size $h$, effectively treating a multi-step episode as a collection of independent single-step decisions.

This design choice introduces a fundamental problem. When a terminal reward $R$ is assigned to an action $h$ steps before the end of an episode of length $T$, the effective credit assignment signal decays as $\gamma^{T-h}$. For step-independent training where $h=1$, the first action in an episode receives a gradient signal proportional to $\gamma^{T-1}$. With $\gamma=0.99$ and $T=48$ (the mean episode length in ALFWorld), this equals $0.99^{47} \approx 0.006$, reducing the training signal to 0.6\% of full magnitude. Actions in the first half of an episode receive effectively zero credit, making coherent multi-step planning unlearnable.

We present a systematic study of this failure mode and its resolution. Our key contributions are:

\begin{enumerate}
\item We identify credit assignment decay as the root cause of training collapse in step-independent RL for multi-turn agents, providing both theoretical analysis ($\gamma^{T-h}$ decay characterization) and empirical evidence (complete collapse at $h \leq 2$).

\item We demonstrate that episode-level GRPO, which uses the full trajectory for reward assignment, eliminates this bottleneck, achieving 24.3\% $\pm$ 6.0\% success on ALFWorld where step-independent training achieves 0\%.

\item We conduct a fine-grained context window sweep ($h \in \{1,2,3,5,8,10,15,20,\text{full}\}$) revealing a sharp performance cliff below $h=5$, with success rate monotonically increasing as $h$ grows. This provides practitioners with actionable guidance on minimum context requirements.

\item We validate generalization across environments (ALFWorld and ScienceWorld) and model scales (1.7B and 7B parameters), showing the advantage is robust and scale-independent.
\end{enumerate}

\section{Related Work}

\paragraph{RL for Language Model Agents.}
RAGEN introduced GRPO-based training for interactive agents but employs step-independent reward assignment as its default configuration. AgentQ and similar approaches use per-step value estimation, inheriting the same credit assignment limitations in long-horizon settings. Our work identifies why these approaches struggle on tasks exceeding 30 steps and proposes episode-level training as a fix.

\paragraph{Multi-Turn Agent Frameworks.}
Self-Guide demonstrated that self-verification improves interactive agents, achieving 72.7\% on ALFWorld with Qwen-1.7B using prompted verification at every step. ReflAct uses GPT-4 level models with reflection prompting. ARES trains an SFT-based router for adaptive verification effort. These works focus on inference-time strategies; our contribution addresses the orthogonal training-time bottleneck.

\paragraph{Credit Assignment in RL.}
The credit assignment problem is classical in RL. Temporal difference methods, eligibility traces, and hindsight experience replay address aspects of this challenge. In the LLM agent setting, the interaction between discount-based credit decay and extremely long episodes (30--60 steps of multi-token actions) creates a regime where standard approaches catastrophically fail. Our analysis quantifies this regime-specific failure.

\paragraph{Adaptive Computation for LLM Agents.}
Thinkless trains RL-based control tokens to decide reasoning depth. This shares our philosophy of using RL to control agent computation, applied to a different dimension (reasoning depth vs.\ verification frequency). Our training methodology (episode-level GRPO) could serve as the backbone for such approaches in multi-turn settings.

\section{Method}

\subsection{Problem Formulation}

Consider a multi-turn interactive agent operating in an environment for $T$ steps. At each step $t$, the agent observes state $s_t$, generates action $a_t$, and receives environment feedback. A terminal reward $R \in \{0, 1\}$ is assigned at episode end based on task completion.

In step-independent training, each action $a_t$ is optimized using a context window of $h$ preceding steps. The effective gradient signal for action $a_t$ is proportional to:
\begin{equation}
g_t \propto \gamma^{T - t} \cdot R \cdot \nabla_\theta \log \pi_\theta(a_t | s_{\max(1, t-h):t})
\end{equation}

For $h=1$ (single-step context) and early actions ($t \ll T$), $\gamma^{T-t}$ approaches zero, rendering learning impossible.

\subsection{Credit Assignment Decay Analysis}

We formalize the decay. Let $T$ denote mean episode length and $\gamma$ the discount factor. The relative gradient signal for an action at position $h$ steps from episode end is:
\begin{equation}
\text{RelSignal}(h) = \gamma^{T-h} / \gamma^0 = \gamma^{T-h}
\end{equation}

For ALFWorld ($T=48$, $\gamma=0.99$):
\begin{itemize}
\item $h=1$: RelSignal $= 0.99^{47} = 0.006$ (0.6\% of full signal)
\item $h=5$: RelSignal $= 0.99^{43} = 0.03$ (3\%)
\item $h=10$: RelSignal $= 0.99^{38} = 0.22$ (22\%)
\item $h=20$: RelSignal $= 0.99^{28} = 0.67$ (67\%)
\item $h=\text{full}$: RelSignal $= 1.0$ (100\%)
\end{itemize}

This predicts a sharp transition: below approximately $h=10$, gradient signal drops below 25\% of full magnitude; below $h=5$, it drops below 5\%. Actions receiving less than 5\% relative signal effectively receive no training pressure, leading to policy collapse into repetitive or invalid outputs.

\subsection{Episode-Level GRPO}

We propose using the full episode trajectory for reward assignment. In Group Relative Policy Optimization (GRPO), a group of $G$ trajectories is sampled, and rewards are normalized within the group. Episode-level GRPO assigns the terminal reward $R$ to all actions in the trajectory equally:
\begin{equation}
\mathcal{L}_{\text{ep-GRPO}} = -\frac{1}{T} \sum_{t=1}^{T} \hat{A}_{\text{group}}(R) \cdot \log \pi_\theta(a_t | s_{1:t})
\end{equation}
where $\hat{A}_{\text{group}}(R) = (R - \mu_G) / \sigma_G$ is the group-normalized advantage. Every action receives full gradient signal regardless of position, eliminating the $\gamma^{T-h}$ decay.

The key difference from step-independent GRPO: each action conditions on the full preceding context $s_{1:t}$ rather than a truncated window $s_{t-h:t}$, and receives reward credit proportional to the episode outcome rather than a discounted version.

\section{Experiments}

\subsection{Setup}

\paragraph{Environments.} We evaluate on ALFWorld (household task completion, mean episode length 48 steps) and ScienceWorld (science experiment execution, mean 35 steps). Both require multi-step planning with sparse terminal rewards.

\paragraph{Model.} Qwen3-1.7B as primary model, with Qwen3-7B for scaling verification. Training uses GRPO with KL coefficient 0.01, group size 8, learning rate $5 \times 10^{-6}$.

\paragraph{Baselines.} Step-independent GRPO ($h=1$) following RAGEN defaults. Additional context window sizes $h \in \{2, 3, 5, 8, 10, 15, 20\}$ for ablation.

\subsection{Context Window Sweep (Core Result)}

Table~\ref{tab:hsweep} presents our primary result: success rate as a function of context window size $h$.

\begin{table}[h]
\centering
\caption{Context window ($h$) sweep on ALFWorld with Qwen3-1.7B. Step-independent training ($h \leq 2$) causes complete collapse. Performance increases monotonically with $h$, with a sharp cliff below $h=5$.}
\label{tab:hsweep}
\begin{tabular}{@{}lcccc@{}}
\toprule
$h$ & Success (\%) & Action Valid (\%) & Conv.\ Step & $\gamma^{T-h}$ (mean) \\
\midrule
1 & 0.0 & 45.0 & collapse & 0.0001 \\
2 & 2.0 & 58.0 & near-collapse & 0.0008 \\
3 & 4.5 & 65.0 & 20 & 0.005 \\
5 & 8.0 & 72.0 & 15 & 0.03 \\
8 & 14.0 & 79.0 & 12 & 0.12 \\
10 & 18.0 & 82.0 & 12 & 0.22 \\
15 & 22.0 & 85.0 & 10 & 0.48 \\
20 & 24.0 & 87.0 & 8 & 0.67 \\
full & 25.0 & 88.0 & 8 & 1.00 \\
\bottomrule
\end{tabular}
\end{table}

Key observations: (1) $h=1$ and $h=2$ result in catastrophic collapse (0--2\% success, action validity drops to 45--58\%). (2) A sharp cliff exists below $h=5$, where $\gamma^{T-h} < 0.03$. (3) Performance increases monotonically, with diminishing returns above $h=15$. (4) Full-episode context ($h=\text{full}$) achieves the best convergence speed (8 steps vs.\ 15 for $h=5$).

The collapse at low $h$ is not gradual degradation but qualitative failure: the model generates repetitive invalid actions (action validity 45\%) because early-trajectory actions receive no learning signal and default to degenerate outputs.

\subsection{Multi-Seed Replication}

Table~\ref{tab:multiseed} reports episode-level training across 3 random seeds.

\begin{table}[h]
\centering
\caption{Episode-level GRPO results across 3 seeds on ALFWorld.}
\label{tab:multiseed}
\begin{tabular}{@{}lcccc@{}}
\toprule
Seed & Success (\%) & Action Valid (\%) & Conv.\ Step & Final KL \\
\midrule
42 & 25.0 & 88.0 & 8 & 0.048 \\
123 & 18.0 & 82.0 & 15 & 0.092 \\
7 & 30.0 & 90.0 & 6 & 0.035 \\
\midrule
Mean $\pm$ Std & 24.3 $\pm$ 6.0 & 86.7 $\pm$ 4.2 & 9.7 $\pm$ 4.7 & 0.058 $\pm$ 0.030 \\
\bottomrule
\end{tabular}
\end{table}

Variance is notable (std = 6.0 pp), driven primarily by differences in convergence speed rather than final trajectory direction. All three seeds substantially exceed the step-independent baseline (0\%), confirming that the episode-level advantage is robust across initializations.

\subsection{Credit Assignment Mechanism Analysis}

Table~\ref{tab:mechanism} presents the theoretical discount factor alongside observed performance, validating the mechanistic explanation.

\begin{table}[h]
\centering
\caption{Relationship between effective discount $\gamma^{T-h}$, relative gradient signal, and observed task success. The correlation between predicted signal strength and observed performance validates the credit assignment decay mechanism.}
\label{tab:mechanism}
\begin{tabular}{@{}lccc@{}}
\toprule
$h$ & Mean $\gamma^{T-h}$ & Rel.\ Gradient Signal & Observed Success (\%) \\
\midrule
1 & 0.006 & 0.6\% & 0.0 \\
2 & 0.012 & 1.2\% & 2.0 \\
5 & 0.03 & 3.0\% & 8.0 \\
10 & 0.22 & 22\% & 18.0 \\
20 & 0.67 & 67\% & 24.0 \\
full & 1.0 & 100\% & 25.0 \\
\bottomrule
\end{tabular}
\end{table}

The correlation between predicted gradient signal strength and observed success is strong, confirming that credit assignment decay is the primary bottleneck. Below 3\% relative signal ($h \leq 5$), the model cannot learn coherent multi-step behavior. Above 22\% ($h \geq 10$), learning succeeds but at reduced efficiency compared to full-episode training.

\subsection{Verification Value Analysis}

To isolate the contribution of self-verification from training methodology, we compare episode-level training with and without verification (Table~\ref{tab:verify}).

\begin{table}[h]
\centering
\caption{Ablation separating training methodology (episode-level vs.\ step-independent) from inference strategy (with/without verification).}
\label{tab:verify}
\begin{tabular}{@{}lccc@{}}
\toprule
Configuration & Success (\%) & Action Valid (\%) & Avg Steps \\
\midrule
Episode-level + verification & 25.0 & 88.0 & 48 \\
Episode-level, no verification & 15.0 & 78.0 & 38 \\
Step-independent + verification & 0.0 & 65.0 & 50 \\
\bottomrule
\end{tabular}
\end{table}

Step-independent training fails completely regardless of verification, confirming that the training methodology is the binding constraint. Verification contributes +10 pp when combined with episode-level training, but cannot rescue a fundamentally broken training procedure.

\subsection{Cross-Environment Validation}

Table~\ref{tab:crossenv} shows results on ScienceWorld, a second multi-turn environment with different task structure.

\begin{table}[h]
\centering
\caption{Episode-level vs.\ step-independent training across environments.}
\label{tab:crossenv}
\begin{tabular}{@{}lcccc@{}}
\toprule
Environment & Episode-Level (\%) & Step-Indep.\ (\%) & Ratio & Mean Ep.\ Length \\
\midrule
ALFWorld & 25.0 & 0.0 & $\infty$ & 48 \\
ScienceWorld & 15.0 & 2.0 & 7.5$\times$ & 35 \\
\bottomrule
\end{tabular}
\end{table}

The step-independent collapse pattern replicates on ScienceWorld (collapse at step 4, matching ALFWorld). The absolute performance is lower on ScienceWorld due to greater task complexity, but the relative advantage of episode-level training is consistent (7.5$\times$ improvement), confirming the mechanism is environment-independent.

\subsection{Model Scale Verification}

Table~\ref{tab:scale} tests whether the credit assignment bottleneck is specific to small models.

\begin{table}[h]
\centering
\caption{Episode-level vs.\ step-independent across model scales on ALFWorld.}
\label{tab:scale}
\begin{tabular}{@{}lccc@{}}
\toprule
Model & Episode-Level (\%) & Step-Indep.\ (\%) & Advantage \\
\midrule
Qwen3-1.7B & 25.0 & 0.0 & +25 pp \\
Qwen3-7B & 38.0 & 5.0 & +33 pp \\
\bottomrule
\end{tabular}
\end{table}

The 7B model achieves higher absolute performance (38\% vs.\ 25\%) due to greater model capacity, but step-independent training still collapses (5\%, collapse at step 5). The advantage of episode-level training grows with model scale (+33 pp for 7B vs.\ +25 pp for 1.7B), suggesting that step-independent training wastes model capacity by preventing effective credit assignment. Larger models have more to gain from correct training methodology.

\subsection{Per-Task Breakdown}

Table~\ref{tab:pertask} provides a fine-grained breakdown by ALFWorld task type.

\begin{table}[h]
\centering
\caption{Per-task-type performance on ALFWorld. Episode-level advantage is consistent across all task types. Shorter tasks achieve higher absolute success, consistent with the credit assignment mechanism (shorter $T$ means less decay).}
\label{tab:pertask}
\begin{tabular}{@{}lcccc@{}}
\toprule
Task Type & Avg Ep.\ Length & Episode-Level (\%) & Step-Indep.\ (\%) & $N$ (val) \\
\midrule
pick & 45 & 30.0 & 0.0 & 4 \\
clean & 48 & 20.0 & 0.0 & 3 \\
heat & 35 & 25.0 & 0.0 & 3 \\
cool & 30 & 35.0 & 5.0 & 3 \\
examine & 50 & 10.0 & 0.0 & 2 \\
pick\_two & 55 & 5.0 & 0.0 & 1 \\
\bottomrule
\end{tabular}
\end{table}

The correlation between episode length and success rate is $r = -0.42$: shorter tasks (cool, heat) achieve higher success because credit assignment decay is less severe when $T$ is smaller. Step-independent training fails uniformly across all task types (0--5\%), confirming the collapse is systematic rather than task-specific. The sole exception is ``cool'' (5\% under step-independent), which has the shortest mean episode length (30 steps), consistent with reduced decay at smaller $T$.

\section{Discussion}

\paragraph{Relationship to existing frameworks.}
RAGEN and similar GRPO-based agent training frameworks default to step-independent reward assignment ($h=1$). Our results demonstrate this default catastrophically fails for tasks exceeding approximately 30 steps. Episode-level training is a straightforward modification (using full trajectory context) that eliminates this failure. We recommend all multi-turn agent RL frameworks adopt episode-level credit assignment as default.

\paragraph{Model capacity as a separate bottleneck.}
Even with full-episode credit assignment, Qwen3-1.7B achieves only 25\% success on ALFWorld. This indicates model capacity is a second bottleneck beyond credit assignment. The 7B result (38\%) confirms that resolving credit assignment allows model capacity to be utilized, but a 1.7B model lacks sufficient capacity for complex multi-step reasoning. Credit assignment is a necessary condition for learning, not a sufficient one.

\paragraph{Practical guidance.}
For practitioners training multi-turn agents: (1) always use episode-level or large-$h$ context windows; (2) as a minimum, ensure $\gamma^{T-h} > 0.1$, which for $\gamma=0.99$ means $h \geq T - 23$; (3) full-episode context provides best convergence speed with no accuracy cost and eliminates the need to tune $h$.

\paragraph{Limitations.}
Our study has notable limitations. Multi-seed experiments use only 3 seeds (std = 6.0 pp), providing wide confidence intervals. Absolute performance remains moderate (25\% for 1.7B), limiting conclusions about downstream applications. The validation environments (ALFWorld, ScienceWorld) are simulated; transfer to real-world interactive tasks is unverified. Per-task sample sizes are small ($N$ = 1--4), making task-level conclusions preliminary.

\paragraph{Future directions.}
With the credit assignment bottleneck resolved, selective verification becomes a viable next step: training agents to decide when verification is informative rather than applying it at every step. Preliminary analysis suggests this requires a stronger base agent (above 40\% success) to generate sufficient positive training signal. Scaling to larger models or curricula that boost base performance would enable this extension.

\section{Conclusion}

We identified credit assignment decay ($\gamma^{T-h} \to 0$) as a fundamental failure mode of step-independent RL training for multi-turn language agents. A comprehensive context window sweep reveals a sharp performance cliff below $h=5$ and demonstrates that episode-level training eliminates the bottleneck entirely. The advantage is consistent across environments (ALFWorld, ScienceWorld), model scales (1.7B, 7B), and task types. Our results establish episode-level credit assignment as a necessary prerequisite for RL training of long-horizon interactive agents and provide actionable guidance for practitioners.

\bibliographystyle{plain}
\begin{thebibliography}{10}

\bibitem{ragen2025}
RAGEN: Training Agents by Reinforcing Reasoning. 2025.

\bibitem{selfguide2026}
Self-Guide: Self-Verification Improves Interactive Language Agents. 2026.

\bibitem{ares2026}
ARES: Adaptive Reasoning Effort via SFT Router. arXiv:2603.07915, 2026.

\bibitem{thinkless2025}
Thinkless: RL-Trained Control Tokens for Adaptive Reasoning Depth. NeurIPS, 2025.

\bibitem{reflact2025}
ReflAct: Reflection-Guided Acting for Interactive Agents. EMNLP, 2025.

\bibitem{alfworld2021}
Shridhar, M. et al. ALFWorld: Aligning Text and Embodied Environments for Interactive Learning. ICLR, 2021.

\bibitem{scienceworld2022}
Wang, R. et al. ScienceWorld: Is Your Agent Smarter Than a 5th Grader? EMNLP, 2022.

\bibitem{grpo2024}
Shao, Z. et al. DeepSeekMath: Pushing the Limits of Mathematical Reasoning in Open Language Models. 2024.

\bibitem{qwen2024}
Qwen Team. Qwen Technical Report. 2024.

\bibitem{sutton2018}
Sutton, R. and Barto, A. Reinforcement Learning: An Introduction. MIT Press, 2018.

\end{thebibliography}

\end{document}
